\chapter{Especificação de Requisitos de Software (SRS)}
\label{chap:especificacao_requisitos}

\section{Visão Geral dos Requisitos}
Esta seção estabelece o escopo e o contexto operacional dos requisitos de software, estruturados em conformidade com as recomendações das normas \textbf{IEEE Std 830-1998} e \textbf{ISO/IEC/IEEE 29148:2018}.

\subsection{Funções do Produto}
Resumo das principais capacidades funcionais oferecidas pelo software aos usuários.

\subsection{Classes e Características dos Usuários}
\begin{itemize}
    \item \textbf{Administrador do Sistema:} Responsável pela gestão de usuários, controle de permissões e auditoria;
    \item \textbf{Operador / Usuário Especialista:} Executa as atividades fins do domínio de aplicação;
    \item \textbf{Cliente / Usuário Final:} Consulta informações e interage via interface pública.
\end{itemize}

\subsection{Restrições Gerais de Design e Implementação}
Restrições tecnológicas, regulatórias, de compatibilidade e segurança impostas à solução.

\section{Requisitos Funcionais (RF)}
Os requisitos funcionais estão catalogados por módulos, identificados de forma unívoca.

\begin{table}[htbp]
\caption{Quadro de Requisitos Funcionais do Sistema}
\label{tab:requisitos_funcionais}
\centering
\small
\begin{tabularx}{\textwidth}{@{} l l Y c @{}}
\toprule
\textbf{ID} & \textbf{Módulo} & \textbf{Descrição Funcional} & \textbf{Prioridade} \\
\midrule
RF01 & Identidade & Permitir autenticação de usuários via JWT com bloqueio após 5 tentativas. & \badgealta \\
RF02 & Agendamento & Registrar agendamento com verificação automática de conflito de agenda. & \badgealta \\
RF03 & Prontuário & Permitir registro de evolução médica com histórico imutável assinado. & \badgealta \\
RF04 & Estoque & Atualizar saldo em tempo real e emitir alerta de ponto de reposição. & \badgemedia \\
RF05 & Relatórios & Gerar relatórios consolidados em formato PDF e planilha eletrônica. & \badgebaixa \\
\bottomrule
\end{tabularx}
\end{table}

\section{Atributos e Dicionário de Dados dos Requisitos}
\begin{table}[htbp]
\caption{Dicionário de Atributos do Requisito de Agendamento (RF02)}
\label{tab:dicionario_dados}
\centering
\small
\begin{tabularx}{\textwidth}{@{} l l l c Y @{}}
\toprule
\textbf{Campo} & \textbf{Tipo} & \textbf{Tamanho / Formato} & \textbf{Obrigatório} & \textbf{Regra de Validação} \\
\midrule
id\_agendamento & UUID & 36 caracteres & Sim & Chave primária gerada pelo sistema. \\
data\_hora & Datetime & ISO 8601 & Sim & Deve ser no futuro e em horário comercial. \\
paciente\_id & UUID & 36 caracteres & Sim & Deve existir e estar ativo no cadastro. \\
status & Enum & [Agendada, Concluída, Cancelada] & Sim & Inicial padrão: Agendada. \\
\bottomrule
\end{tabularx}
\end{table}

\section{Restrições de Domínio e Regras de Negócio (RN)}
As regras de negócio do sistema regem as validações fundamentais de integridade operacional:

\begin{table}[htbp]
\caption{Catálogo de Regras de Negócio (RN)}
\label{tab:regras_negocio}
\centering
\small
\renewcommand{\arraystretch}{1.25}
\begin{tabularx}{\textwidth}{@{} l l Y @{}}
\toprule
\textbf{Código} & \textbf{Nome da Regra} & \textbf{Descrição e Condição de Guarda} \\
\midrule
RN01 & Validação de Conflito de Horários & Não é permitido que o mesmo profissional tenha dois atendimentos agendados com sobreposição temporal no mesmo consultório ou horário. \\
RN02 & Imutabilidade de Registros Clínicos & Registros de atendimento e prescrições médicas finalizados não podem ser alterados nem excluídos, sendo permitidas apenas retificações mediante novas entradas assinadas digitalmente. \\
\bottomrule
\end{tabularx}
\end{table}

\section{Requisitos Não Funcionais (RNF)}
\begin{table}[htbp]
\caption{Quadro de Requisitos Não Funcionais}
\label{tab:requisitos_nao_funcionais}
\centering
\small
\begin{tabularx}{\textwidth}{@{} l l Y l @{}}
\toprule
\textbf{ID} & \textbf{Categoria} & \textbf{Especificação e Métrica} & \textbf{Padrão / Critério} \\
\midrule
RNF01 & Desempenho & O tempo de resposta para operações comuns deve ser $\le 2{,}0$ segundos sob carga nominal de 100 requisições/s. & $95\%$ das requisições. \\
RNF02 & Segurança & Senhas devem ser armazenadas com hash Argon2 ou Bcrypt com salt aleatório. & OWASP Top 10. \\
RNF03 & Disponibilidade & O sistema deve manter índice de disponibilidade de no mínimo $99{,}5\%$ em produção. & SLA acordado. \\
RNF04 & Conformidade & Todos os dados pessoais sensíveis devem ser tratados em conformidade estrita com a LGPD. & Lei nº 13.709/2018. \\
\bottomrule
\end{tabularx}
\end{table}
